\section{结论}
针对问题一，本文建立并求解了圆柱径向传热与非线性水分扩散模型，给出七个指定时刻的温度、干基含水率分布，以及符合模板工作表名与距离列结构的完整结果文件。预热阶段呈现表面先升温、表层先失水、中心响应滞后的特征。模型的离散守恒与数值收敛已通过验证，结论适用于本文明确列出的有效边界及忽略蒸发冷却等假设。
